% !TeX root = ../../Book3.tex
% 纲要标题：### 13.2 闭值域与 Fredholm 择一定理
% 纲要位置：几何学计划纲要.md，第 2619 行。
% 正文状态：初稿；尚未正式审核。

% 纲要内容开始
% 
% * 值域的闭性
% * 形式伴随的正则性
% * 余核与伴随核
% * 可解条件和 Fredholm 指数
% * 指数对连续扰动的稳定性
% 
% 纲要内容结束

\section{闭值域与 Fredholm 择一定理}
\label{g3:sec:1302}

有了核上的投影与补空间估计，方程能否求解就转化为右端是否落在值域中。闭值域保证近似可解的右端确实可解；形式伴随把所有障碍写成有限个积分条件。这是本节的 Fredholm 选择。\footnote{Erik Ivar Fredholm（1866--1927），瑞典数学家，研究积分方程、谱理论与数学物理。}它把矩阵方程中“右端与转置核正交”的条件推广到闭流形上的椭圆方程。

\subsection{值域的闭性}
\label{g3:subsec:1302:01}
若 \(Pu_j\) 收敛，\(u_j\) 未必收敛，因为可以给它任意加上核元素。先减去核投影，上一节的估计就把右端的 Cauchy 性传给解。

\begin{proposition}\label{g3:prop:elliptic-closed-range}
	设 \(M^n\) 为非空闭光滑流形，\(n\ge1\)，\(E,F\) 为同秩光滑有限秩度量丛，\(P:C^\infty(E)\to C^\infty(F)\) 是光滑椭圆算子，阶 \(d\ge1\)。固定底流形的光滑黎曼度量以及各丛的光滑参考联络，Sobolev范数用这些数据定义；积分与形式伴随采用该度量的体积密度及给定的纤维度量。记 \(\Pi_P\) 为到光滑核 \(\ker P\) 的 \(L^2\) 正交投影。则 \(P:H^d(E)\to L^2(F)\) 的值域 \(\mathcal R_P\) 闭，且
	\[
		P:\{u\in H^d(E):\Pi_Pu=0\}\longrightarrow\mathcal R_P
	\]
	为双射，其逆有界。
\end{proposition}
\begin{proof}
	设 \(f_j=Pu_j\to f\) 于 \(L^2\)，令 \(v_j=u_j-\Pi_Pu_j\)。\Cref{g3:thm:kernel-orthogonal-estimate}给出
	\[
		\|v_j-v_\ell\|_{H^d}\le C\|f_j-f_\ell\|_2.
	\]
	故 \(v_j\) 在完备空间 \(H^d\) 中收敛到 \(v\)。算子的有界性给出 \(Pv=f\)，投影的连续性给出 \(\Pi_Pv=0\)，所以 \(f\in\mathcal R_P\)。核正交截面中的齐次解为零，故所述映射单射；减去核投影又不改变右端，故它满射。逆的范数界正是上一节的估计。
\end{proof}

证明中的关键下界来自上一节的矛盾紧性论证。闭值域使残差趋零具有求解意义：只要右端 \(f\) 可以被某列 \(Pu_j\) 在 \(L^2\) 中逼近，就能在去掉核以后得到一个真正满足 \(Pu=f\) 的极限。

\begin{example}\label{g3:exm:dense-nonclosed-range}
	在 \(\ell^2\) 中令 \(T(x_j)=(x_j/j)\)。它有界且单射，值域包含所有有限支序列，因而稠密。令 \(y=(1/j)_{j\ge1}\in\ell^2\)，其有限截断均在值域中，且收敛到 \(y\)；但 \(Tx=y\) 要求 \(x_j=1\)，不属于 \(\ell^2\)，故值域不闭。这里伴随核为零，却仍有不可解右端，说明正交条件要成为充分条件，闭值域不能缺少。
\end{example}

\subsection{形式伴随的正则性}
\label{g3:subsec:1302:02}
第十二章的形式伴随 \(P^*\) 由积分恒等式确定：
\[
	(Pu,v)_2=(u,P^*v)_2,\qquad u\in C^\infty(E),\quad v\in C^\infty(F).
\]
闭流形没有边界项。按\cref{g3:prop:symbol-calculus}，其主符号为 \(\sigma_d(P^*)(x,\xi)=\sigma_d(P)(x,\xi)^*\)，故 \(P^*\) 同样椭圆，阶数仍为 \(d\)。于是 \(P^*v=0\) 的分布解也是光滑的，且 \(\mathcal N_{P^*}\) 有限维。

\ProseParagraphBreak

形式伴随是微分表达式。若把 \(P\) 看成从 \(H^d\) 到 \(L^2\) 的有界映射，它关于这两个 Hilbert 范数的有界伴随取值于 \(H^d\)，并不是同一个表达式。谱论需要的是把 \(P\) 看成 \(L^2\) 上定义域为 \(H^d\) 的算子，下面把这一点准确落实。

\begin{proposition}\label{g3:prop:elliptic-adjoint-domain}
	设 \(M^n\) 为非空闭光滑流形，\(n\ge1\)，\(E,F\) 为同秩有限秩光滑 Hermitian 丛，\(P:C^\infty(E)\to C^\infty(F)\) 为光滑椭圆 \(d\) 阶算子，\(d\ge1\)。固定底流形的光滑黎曼度量以及各丛的光滑参考联络，Sobolev范数用这些数据定义；积分与形式伴随采用该度量的体积密度及给定的纤维度量。以 \(P^*\) 表示由固定体积密度及丛度量确定的形式伴随。把 \(\mathsf P\) 定义为 \(L^2(E)\to L^2(F)\) 的算子，域为 \(H^d(E)\)，作用为 \(P\)。则 \(\mathsf P\) 稠定且闭。其 Hilbert 伴随的域恰为 \(H^d(F)\)，在此域上作用为 \(P^*\)。
\end{proposition}
\begin{proof}
	光滑截面在 \(L^2\) 中稠密，故定义域稠密。若 \(u_j\to u\) 于 \(L^2(E)\)、\(Pu_j\to f\) 于 \(L^2(F)\)，全局估计给出
	\[
		\|u_j-u_\ell\|_{H^d}
		\le C\bigl(\|Pu_j-Pu_\ell\|_2+\|u_j-u_\ell\|_2\bigr).
	\]
	因此 \(u_j\) 在 \(H^d\) 中收敛，极限在 \(L^2\) 中只能是 \(u\)，且 \(Pu=f\)。这正是图闭的含义。

	按定义，\(v\in L^2(F)\) 属于 Hilbert 伴随的域，等价于存在 \(w\in L^2(E)\)，使
	\((Pu,v)_2=(u,w)_2\) 对全部 \(u\in H^d(E)\) 成立。对光滑 \(u\) 作测试，得 \(P^*v=w\) 分布成立。伴随的椭圆正则性使 \(v\in H^d(F)\)。反之，若 \(v\in H^d(F)\)，先对光滑 \(u,v\) 分部积分，再由 \(H^d\) 光滑逼近取极限，得到上述恒等式，其中 \(w=P^*v\in L^2\)。故两种定义域恰好相同。
\end{proof}

\subsection{余核与伴随核}
\label{g3:subsec:1302:03}
对右端空间中的闭子空间 \(\mathcal R_P\)，正交补中的向量记录了全部线性阻碍。形式伴随的正则性把这些原本只有 \(L^2\) 意义的向量变成光滑截面。

\termidx[Fredholm xuanze]{Fredholm 选择}
\foreignidx{Fredholm alternative}
\begin{theorem}[椭圆 Fredholm 选择]\label{g3:thm:elliptic-fredholm}
	设 \(M^n\) 为非空闭光滑流形，\(n\ge1\)，\(E,F\) 为同秩光滑有限秩度量丛，\(P:C^\infty(E)\to C^\infty(F)\) 为光滑椭圆算子，阶 \(d\ge1\)。固定底流形的光滑黎曼度量以及各丛的光滑参考联络，Sobolev范数用这些数据定义；积分与形式伴随采用该度量的体积密度及给定的纤维度量。记 \(P^*\) 为形式伴随，\(\mathcal N_Q\) 为 \(Q=P,P^*\) 在光滑截面上的核，\(\Pi_P\) 为到 \(\mathcal N_P\) 的 \(L^2\) 正交投影。则
	\begin{equation}\label{g3:eq:elliptic-range-orthogonal}
		\operatorname{ran}(P:H^d(E)\to L^2(F))
		=\mathcal N_{P^*}^{\perp}.
	\end{equation}
	因此 \(L^2(F)/\operatorname{ran}P\) 自然同构于 \(\mathcal N_{P^*}\)，特别地余维有限。对 \(f\in L^2(F)\)，方程 \(Pu=f\) 可解当且仅当
	\[
		(f,h)_2=0\qquad\text{对全部 }h\in\mathcal N_{P^*}.
	\]
	可解时有唯一满足 \(\Pi_Pu=0\) 的 \(H^d\) 解，其余解为它加上任意核元素。
\end{theorem}
\begin{proof}
	若 \(v\in\mathcal R_P^\perp\)，则 \((Pu,v)_2=0\) 对全部光滑 \(u\) 成立，所以 \(P^*v=0\) 分布成立。正则性给出 \(v\in\mathcal N_{P^*}\)。反向由分部积分和光滑稠密性立得，故
	\(\mathcal R_P^\perp=\mathcal N_{P^*}\)。\Cref{g3:prop:elliptic-closed-range}保证值域闭；Hilbert 投影定理\cref{g3:thm:hilbert-projection}因而给出
	\(\mathcal R_P=(\mathcal R_P^\perp)^\perp\)，得到式\eqref{g3:eq:elliptic-range-orthogonal}。

	余核的具体同构把类 \([f]\) 送到 \(\Pi_{P^*}f\)：给 \(f\) 加上值域元素不改变投影，投影为零又恰好意味着 \(f\) 在值域中，故它良定义且双射。存在性随值域等式成立；减去 \(\Pi_Pu\) 给出归一化解，两个这样的解之差同时属于核和核正交补，只能为零。
\end{proof}

\begin{corollary}\label{g3:cor:elliptic-range-hk}
	设 \(M\) 为非空闭光滑黎曼流形，\(E,F\to M\) 为带固定底流形的光滑黎曼度量，以及两丛的光滑纤维度量和参考联络的同秩有限秩实或复光滑向量丛，\(P:C^\infty(E)\to C^\infty(F)\) 为 \(d\ge1\) 阶光滑椭圆算子。记 \(P^*\) 为形式伴随，\(\mathcal N_{P^*}=\ker(P^*:C^\infty(F)\to C^\infty(E))\)，对每个整数 \(k\ge0\)，有
	\[
		P(H^{k+d}(E))=\{f\in H^k(F):(f,h)_2=0\text{ 对所有 }h\in\mathcal N_{P^*}\}.
	\]
	相容右端的核正交解满足 \(\|u\|_{H^{k+d}}\le C_k\|f\|_{H^k}\)。特别地，光滑右端的解光滑。
\end{corollary}
\begin{proof}
	必要性由分部积分。相容的 \(H^k\) 右端先由上一结果得到 \(H^d\) 解，再由\cref{g3:thm:global-elliptic-hk}提升到 \(H^{k+d}\)。归一化与估计来自\cref{g3:thm:kernel-orthogonal-estimate}；光滑性可由逐阶提升或直接由分布椭圆光滑性得到。
\end{proof}

这里的余核用光滑 \(\mathcal N_{P^*}\) 表示，与 Sobolev 阶数的选择相容。不过在 \(H^k\) 中，所写条件仍是 \(L^2\) 配对，不能自行替换为 \(H^k\) 正交条件。

\subsection{可解条件和 Fredholm 指数}
\label{g3:subsec:1302:04}
\begin{definition}\label{g3:def:fredholm-operator}
	设 \(X,Y\) 为实或复 Banach 空间。若有界线性映射 \(T:X\to Y\) 的核有限维、值域闭且余核 \(Y/\operatorname{ran}T\) 有限维，则称它为\term[Fredholm suanzi]{Fredholm 算子}[Fredholm operator]，定义其\term[Fredholm zhishu]{Fredholm 指数}[Fredholm index]为
	\[
		\operatorname{ind}T=\dim\ker T-\dim(Y/\operatorname{ran}T).
	\]
\end{definition}
\symidx{\(\operatorname{ind}P\)：核维数减余核维数的 Fredholm 指数}

椭圆算子的上述实现因而是 Fredholm 的，指数为
\(\dim\mathcal N_P-\dim\mathcal N_{P^*}\)。这是一个有限整数；它不要求 \(P\) 自伴，也不要求指数为零。若 \(E=F\)、\(P=P^*\)，指数才立即为零。抽象紧算子的 Fredholm 选择及定义可对照 Brezis 第6.2节和第6章评注第1项\cite[第6.2节，第6章评注第1项]{Brezis2011}。

\begin{theorem}\label{g3:thm:elliptic-fredholm-wkp}
	设 \(M^n\) 为非空闭光滑流形，\(n\ge1\)，\(E,F\) 为同秩光滑有限秩度量丛，\(P:C^\infty(E)\to C^\infty(F)\) 为光滑椭圆 \(d\) 阶算子，\(d\ge1\)。固定底流形的光滑黎曼度量 \(g\)以及各丛的光滑参考联络，Sobolev范数用这些数据定义；积分与形式伴随采用该度量的体积密度及给定的纤维度量。记 \(\mathcal N_P,\mathcal N_{P^*}\) 为 \(P\) 及其形式伴随的光滑核，\(\Pi_P\) 为到 \(\mathcal N_P\) 的 \(L^2\) 正交投影。对整数 \(k\ge0\)、\(1<p<\infty\)，有界映射
	\[
		P:W^{k+d,p}(E)\longrightarrow W^{k,p}(F)
	\]
	为 Fredholm 算子，核为 \(\mathcal N_P\)，值域为
	\[
		\{f\in W^{k,p}(F):\textstyle\int_M\langle f,h\rangle\mu_g=0
		\text{ 对全部 }h\in\mathcal N_{P^*}\}.
	\]
	余核由 \(\mathcal N_{P^*}\) 表示，指数独立于 \(k,p\)。相容右端有唯一的 \(\Pi_Pu=0\) 解，并满足 \(\|u\|_{W^{k+d,p}}\le C_{k,p}\|Pu\|_{W^{k,p}}\)。
\end{theorem}
\begin{proof}
	先取 \(k=0\)。上一节的 \(L^p\) 核正交估计使值域闭：对收敛列 \(Pu_j\) 减去核投影后，\(u_j-\Pi_Pu_j\) 在 \(W^{d,p}\) 中为 Cauchy 列，与闭值域命题的证明完全相同。

	令 \(p'=p/(p-1)\)。由积分对偶\cref{g3:lem:lp-duality-reflexivity}及纤维度量，每个 \(L^p(F)\) 上的连续线性泛函都写成
	\(f\mapsto\int_M\langle f,h\rangle\mu_g\)，其中 \(h\in L^{p'}(F)\)。它消灭值域，当且仅当 \(P^*h=0\) 分布成立；于是 \(h\) 光滑且属于 \(\mathcal N_{P^*}\)。反之每个这样的 \(h\) 都消灭值域。若 \(f\) 不在闭值域中，Hahn--Banach 分离\cref{g3:lem:bounded-functional-extension}给出一个在值域上为零而在 \(f\) 上非零的泛函，故必有上述某个积分不为零。这就证明了 \(k=0\) 的值域等式。

	取 \(\mathcal N_{P^*}\) 的 \(L^2\) 正交归一基，积分配对构成满射 \(L^p(F)\to\mathbb K^{\dim\mathcal N_{P^*}}\)：将右端取为该光滑基的线性组合即可得到任意坐标。因此值域的余维恰为这个维数。

	对于 \(k>0\)，相容的 \(W^{k,p}\) 右端先有 \(W^{d,p}\) 解，再由\cref{g3:thm:global-elliptic-wkp}提升正则性。上述有限个积分在 \(W^{k,p}\) 中仍连续且联合满射，所以值域闭、余维相同。核光滑且不随空间改变，唯一性、归一化和估计均由上一节得到。
\end{proof}

\begin{example}\label{g3:exm:circle-fredholm-compatibility}
	在 \(S^1=\mathbb R/(2\pi\mathbb Z)\) 上，\(P=\partial_x\) 的形式伴随为 \(-\partial_x\)，两者的核都由常数生成，指数为零。方程 \(u'=f\) 有周期解当且仅当 \(\int_0^{2\pi}f=0\)；取
	\[
		u(x)=\int_0^x f(t)\,\mathrm dt-
		\frac1{2\pi}\int_0^{2\pi}\int_0^y f(t)\,\mathrm dt\,\mathrm dy
	\]
	即得唯一零均值解。这里可解性并不要求算子自伴。对 \(\Delta_g\) 而言，核在每个连通分支上由常数构成，因此相容条件是逐个分支的积分为零，不能只检查总积分。
\end{example}

\subsection{指数对连续扰动的稳定性}
\label{g3:subsec:1302:05}
核与余核的维数可能同时变化，差却保持不变。把无限维可逆部分先消去，会留下一个有限维映射；它的秩怎样变化都不改变“核维数减余维数”。

\begin{theorem}\label{g3:thm:index-local-constancy}
	设 \(M^n\) 为非空闭光滑流形，\(n\ge1\)，\(E,F\) 为同秩光滑有限秩度量丛，\(P:C^\infty(E)\to C^\infty(F)\) 为光滑椭圆 \(d\) 阶算子，\(d\ge1\)。固定底流形的光滑黎曼度量以及各丛的光滑参考联络，Sobolev范数用这些数据定义；积分与形式伴随采用该度量的体积密度及给定的纤维度量。把它看成 \(X=H^d(E)\to Y=L^2(F)\) 的有界映射。则存在 \(\delta>0\)，使每个有界线性映射 \(Q:X\to Y\) 满足 \(\|Q-P\|<\delta\) 时，\(Q\) 为 Fredholm 算子且
	\[
		\operatorname{ind}Q=\operatorname{ind}P.
	\]
	特别地，在固定 \(M,E,F\) 上，阶数为 \(d\) 的光滑椭圆算子组成的算子范数连续道路 \(P_t:H^d(E)\to L^2(F)\)、\(0\le t\le1\)，具有恒定指数。
\end{theorem}
\begin{proof}
	记 \(K=\mathcal N_P\)、\(C=\mathcal N_{P^*}\)、\(R=\mathcal R_P\)，并令 \(X_0=\ker\Pi_P\cap H^d(E)\)。我们已经证明有连续直和
	\[
		X=X_0\oplus K,\qquad Y=R\oplus C,
	\]
	以及有界可逆映射 \(A=P|_{X_0}:X_0\to R\)。按这些分解把 \(Q\) 写成
	\[
		Q=\begin{pmatrix}A_Q&B_Q\\ C_Q&D_Q\end{pmatrix}.
	\]
	投影均有界，所以 \(Q\) 充分接近 \(P\) 时，\(\|A^{-1}(A_Q-A)\|<1\)。Neumann 级数给出 \(A_Q\) 的有界逆。作有界可逆的行、列消元，得到
	\begin{equation}\label{g3:eq:fredholm-block-reduction}
		\begin{pmatrix}I&0\\-C_QA_Q^{-1}&I\end{pmatrix}
		Q
		\begin{pmatrix}I&-A_Q^{-1}B_Q\\0&I\end{pmatrix}
		=\begin{pmatrix}A_Q&0\\0&S_Q\end{pmatrix},
		\quad S_Q=D_Q-C_QA_Q^{-1}B_Q:K\to C.
	\end{equation}
	右端的第一块可逆，第二块作用在有限维空间之间，故值域为闭空间 \(R\oplus\operatorname{ran}S_Q\)，核与余核均有限维。可逆行、列变换保持这些性质和维数，因而
	\[
		\operatorname{ind}Q
		=\dim\ker S_Q-\bigl(\dim C-\operatorname{rank}S_Q\bigr)
		=\dim K-\dim C.
	\]
	这就是 \(\operatorname{ind}P\)。沿连续道路，每一点附近都可应用同一论证，故指数局部常值；区间连通使它处处相同。
\end{proof}

分块消元把所有无限维求解集中到同构 \(A_Q\)，剩下的有限维映射 \(S_Q\) 决定核与余核怎样跳跃，如图\cref{g3:fig:fredholm-reduction}。其秩同时从源维数与靶维数中扣除，因此在指数里抵消。
\begin{figure}[htbp]
	\centering
	% !TeX root = ../Book3.tex
\begin{tikzcd}[row sep=1.1em,column sep=4em]
	X_0 \arrow[r,"A_Q", "\simeq"'] & R \\
	K \arrow[r,"S_Q"'] & C
\end{tikzcd}

	\caption[可逆部分与有限维剩余方程]{可逆坐标变换后，两块独立作用。上方是闭补空间之间的同构，下方保留有限维核与余核的变化。}
	\label{g3:fig:fredholm-reduction}
\end{figure}

\begin{corollary}\label{g3:cor:lower-order-index-invariance}
	设 \(M\) 为闭光滑流形，\(E,F\) 为同秩光滑有限秩度量丛，\(P:C^\infty(E)\to C^\infty(F)\) 是光滑椭圆 \(d\) 阶算子，\(d\ge1\)。若 \(B:C^\infty(M;E)\to C^\infty(M;F)\) 为光滑系数、阶数至多 \(d-1\) 的算子，则 \(\operatorname{ind}(P+B)=\operatorname{ind}P\)。
\end{corollary}
\begin{proof}
	\(P_t=P+tB\) 的主符号与 \(P\) 相同，所以整条道路椭圆；又有 \(\|P_t-P_s\|_{H^d\to L^2}\le |t-s|\|B\|\)。应用上一结果即可。
\end{proof}

\begin{example}\label{g3:exm:kernel-jump-circle}
	在复值周期函数上取 \(P_t=-\mathrm i\partial_x+t\)，\(|t|<1/2\)。齐次解为 \(u(x)=a e^{-\mathrm itx}\)，它周期当且仅当 \(t=0\) 或 \(a=0\)。算子形式自伴，故核与余核在 \(t=0\) 同时从零维跳到一维，指数始终为零。\(t\ne0\) 时对常数右端一有 \(P_t(1/t)=1\)，因此逆算子的范数至少按 \(1/|t|\) 增长。这解释了为什么指数稳定并不保证解对参数一致有界。
\end{example}
% 纲要条目排为编号小节标题，并进入目录与书签；说明段仍保留于纲要注释中。
